%%%%%%%%%%%%%%%%%%%%%%%%%%%%%%%%%%%%%%%%%%%%%%%%%%%%%%%%%%%%%%%%%%%%%%%%%%%
% Lehreinheit 2: Die Forschungsfrage
% Quelle: WuZ-Buch, Kap. 15 (LE 2), Kap. 4 (Neurokognition, Information Gap),
%         Kap. 7d (Das wohldefinierte Problem), Kap. 11 (Effekttabelle:
%         "Das geheime Objekt")
% Staunmoment: Das geheime Objekt  |  Kernbegriff: Operationalisierung
%%%%%%%%%%%%%%%%%%%%%%%%%%%%%%%%%%%%%%%%%%%%%%%%%%%%%%%%%%%%%%%%%%%%%%%%%%%
\nsl{In der ersten Einheit stand das Staunen selbst im Mittelpunkt. Jetzt
geht es darum, was man mit ihm anfängt. Ein Staunerlebnis erzeugt eine
Frage, aber meist eine unscharfe: \enquote{Wie kann das sein?} Damit lässt
sich keine Untersuchung planen. Diese Einheit zeigt, wie aus einer echten,
oft vagen Frage eine Forschungsfrage wird, die man beantworten kann. Der
Staunmoment ist diesmal kein Kunststück mit geheimer Methode, sondern ein
Satz mit einer verborgenen Voraussetzung. An ihm lässt sich zeigen, dass
jede Frage Annahmen mitbringt, die man prüfen muss, bevor man nach der
Antwort sucht. Die zweite Hälfte der Stunde gehört den Fragen, die die
Studierenden aus LE~1 mitgebracht haben.
\newpage}
\os{\begin{description}
\item[{\rot\bf Heute}]~\\[-\lblskip]
\bi
  \ite ein Objekt, das noch niemand gesehen hat
  \ite vom Thema zur Forschungsfrage
  \ite Neugier als Informationslücke
  \ite Kriterien guter Forschungsfragen: SMART und FINER
  \ite Operationalisierung: vom Begriff zur Messung
  \ite Werkstatt: Ihre eigene Frage
\ei
\end{description}
\newpage}

\begin{vorbereitung}[Das geheime Objekt]
\textbf{Eine Woche vorher}
\begin{checkliste}
  \item \textbf{Allergien abfragen.} In der Mail zur Einheit (oder am Ende
    von LE~1) fragen, ob jemand eine Erdnussallergie hat, ohne zu
    verraten, wozu. Bei schweren Erdnussallergien können schon Schalenstaub
    und Hautkontakt Reaktionen auslösen. Gibt es auch nur eine Meldung,
    wird die Variante mit der Erbsenschote gespielt (siehe unten). Wer
    nicht sicher abfragen kann, nimmt gleich die Erbse.
  \item Einen \textbf{undurchsichtigen Becher} besorgen (Keramikbecher oder
    ein umgedrehter Pappbecher). Er sollte groß genug sein, dass die Form
    des Objekts nicht zu erraten ist, und auf einem Teller oder Brettchen
    stehen, damit man ihn durch die Reihen tragen kann.
  \item \textbf{Erdnüsse in der Schale} (ungesalzen, Supermarkt, eine kleine
    Tüte reicht). Eine Nuss mit unbeschädigter, geschlossener Schale
    aussuchen, zwei Ersatznüsse einpacken. Die Nuss bleibt bis zur
    Auflösung in einem verschließbaren Gefrierbeutel.
  \item \textbf{Variante ohne Allergierisiko:} eine frische, geschlossene
    Erbsenschote (Gemüseabteilung, Wochenmarkt); außerhalb der Saison eine
    Bohne in der Hülse oder eine Esskastanie. Die Pointe funktioniert mit
    jedem Samen in einer geschlossenen Hülle, deren Inneres noch niemand
    gesehen hat.
  \item Sprechtext üben, vor allem die Pausen. Der Satz muss wortgleich
    wiederholt werden können, weil die Studierenden später mit ihm
    arbeiten.
\end{checkliste}

\textbf{Handout und Übungsblatt}
\begin{checkliste}
  \item Handout \enquote{Die Forschungsfrage} und Übungsblatt LE~2 je
    einmal pro Person kopieren (Druckvorlage \texttt{handouts.pdf}, Seiten
    zu LE~2). Das Übungsblatt enthält die Liste mit acht Beispielfragen für
    die Übung zu impliziten Prämissen; es wird erst nach dem Staunmoment
    ausgeteilt.
  \item Den Satz der Behauptung groß auf ein A3-Blatt schreiben oder als
    Folie bereithalten (\enquote{Unter diesem Becher liegt etwas, das noch
    kein Mensch je gesehen hat. Und nach dieser Stunde wird es auch niemand
    mehr sehen.}).
  \item Die Studierenden in der Erinnerungsmail bitten, ihre Frage aus LE~1
    schriftlich mitzubringen. Für Personen ohne eigene Frage stehen
    Ersatzthemen auf dem Übungsblatt (Aufgabe~2).
  \item Einen Stapel Karteikarten (A6) für die Sicherung am Ende.
\end{checkliste}

\textbf{Am Tag selbst}
\begin{checkliste}
  \item Becher mit Objekt \emph{vor} Beginn sichtbar auf das Pult stellen,
    ohne Kommentar. Fragt schon vorher jemand, was darunter ist: nicht
    antworten.
  \item Tafel zweiteilen: links \enquote{Was könnte es sein?}, rechts
    \enquote{Was muss wahr sein, damit der Satz stimmt?} (rechte Hälfte
    zunächst leer lassen).
  \item Zweiergruppen für die Werkstatt so einteilen, dass möglichst
    Personen aus verschiedenen Fächern zusammensitzen; fremde Augen sehen
    Prämissen schneller.
\end{checkliste}

\textbf{Risiken und Plan B.} Wer das Rätsel kennt, löst es sofort. Dann
die Person bitten, die Lösung aufzuschreiben und für sich zu behalten, und
sie in der Auflösung als Zeugin aufrufen: Was hat ihr geholfen, die
Prämisse zu sehen? Rät niemand die Lösung, ist das kein Problem; der
Staunmoment liegt in der Auflösung, nicht im Erraten. Nach der Stunde wird
die Nuss weder im Raum gegessen noch verteilt; sie wandert in den
Gefrierbeutel und wird draußen gegessen oder entsorgt. An der Pointe
ändert das nichts: Ihr Inneres sieht danach niemand mehr.
\end{vorbereitung}

\begin{ablauf}
\begin{description}
\item[0--5 · Blitzlicht.] Zwei bis drei Tagebuch-Staunmomente der Vorwoche
  teilen und sofort in Frageform bringen: \enquote{Was genau möchten Sie
  daran wissen?}
\item[5--15 · Staunmoment.] \enquote{Das geheime Objekt}. Leitfrage:
  Welche unausgesprochene Voraussetzung steckt in der Behauptung? Erst
  Vermutungen sammeln, dann auflösen.
\item[15--45 · Theorie-Input.] Vom Thema zur Frage; Neugier als
  Informationslücke (Kap.~4 des Buchs); Kriterien guter Forschungsfragen
  (SMART, FINER); implizite Prämissen; Operationalisierung;
  Forschungslücke. Gegenüberstellung staunenswerte vs.\ beantwortbare
  Frage.
\item[45--75 · Werkstatt am eigenen Thema.] Zweiergruppen sezieren die
  mitgebrachte Frage: implizite Prämissen freilegen, messbare Größen
  benennen, nach SMART schärfen (Übungsblatt, Aufgaben 1 und 2).
\item[75--90 · Sicherung.] Jede Person schreibt eine erste eigene
  Forschungsfrage auf eine Karteikarte; drei bis vier Karten werden im
  Plenum auf Plausibilität geprüft; Auftrag zur schriftlichen Fassung.
\end{description}
\end{ablauf}

%==============================================================================
\section{Staunmoment: Das geheime Objekt}
%==============================================================================
\begin{staunbox}{Das geheime Objekt}
Die Lehrperson hebt einen geschlossenen Becher hoch.

\medskip
\emph{\enquote{Unter diesem Becher liegt etwas, das noch kein Mensch je
gesehen hat. Und nach dieser Stunde wird es auch niemand mehr sehen.}}

\medskip
Was könnte darunter liegen? Und wie könnte man prüfen, ob der Satz
stimmt?
\end{staunbox}
\nsl{Der Effekt stammt aus der Effekttabelle des Buchs (Kap.~11), wo er
für Biologie, Philosophie und alle diskursiven Fächer vorgeschlagen wird.
Die Lehrperson sagt den Satz langsam und wiederholt ihn auf Nachfrage
wörtlich. Fragen zum Objekt beantwortet sie in dieser Phase nicht, nimmt
aber jede Vermutung auf und schreibt sie an die linke Tafelhälfte.}
\os{\newpage}

\begin{description}
\item[{\rot\bf Was liegt unter dem Becher?}]~\\[-\lblskip]
\bi
  \ite sammeln Sie Vermutungen, auch unwahrscheinliche
  \ite Welche Wörter im Satz lenken Ihre Vermutungen?
  \ite Was \emph{muss} wahr sein, damit der Satz stimmt?
\ei
\nsl{Die Vermutungen gehen meist in Richtung des Seltenen und Wertvollen:
ein Meteoritensplitter, ein versiegelter Brief, eine Zeichnung, etwas
Verbotenes. Sie verraten, wie die Gruppe den Satz verstanden hat: als
Aussage über eine \emph{Art} von Ding, die noch niemand kennt. Wenn die
Liste an der Tafel lang genug ist, wird die rechte Tafelhälfte eröffnet.}
\os{\newpage}

\item[{\rot\bf Auflösung}]~\\[-\lblskip]
\bi
  \ite unter dem Becher: \textbf{eine Erdnuss in ihrer Schale}
  \ite die Schale wird geöffnet: diese Nuss hat noch nie jemand gesehen
  \ite nach der Stunde ist sie gegessen, niemand wird sie wieder sehen
  \ite der Satz war die ganze Zeit wahr
\ei
\nsl{Bei der Variante ohne Allergierisiko öffnet die Lehrperson eine
Erbsenschote. Botanisch ist das sogar die genauere Wahl: Die Erdnuss ist
keine Nuss, sondern eine Hülsenfrucht wie die Erbse. Auch das lässt sich
nutzen, denn schon der Name \enquote{Erdnuss} enthält eine Annahme, die
sich bei näherem Hinsehen als falsch erweist.}
\os{\newpage}

\item[{\rot\bf Wo lag der Trick?}]~\\[-\lblskip]
\bi
  \ite \enquote{etwas, das noch kein Mensch gesehen hat}: wir hören eine
    \emph{Art} von Ding, die unbekannt ist
  \ite gemeint war \emph{dieses eine} Exemplar
  \ite die verborgene Prämisse lag nicht im Becher, sondern in unserer
    Lesart des Satzes
  \ite jede Frage, jede Behauptung bringt solche Prämissen mit
\ei
\os{\newpage}
\end{description}

%==============================================================================
\section{Vom Thema zur Forschungsfrage}
%==============================================================================
\nsl{Wer eine Arbeit beginnt, hat in der Regel zunächst ein Thema:
\enquote{Zauberkunst im Unterricht}, \enquote{Homeoffice}, \enquote{Social
Media und Jugendliche}. Ein Thema ist ein Gebiet, keine Aufgabe. Man kann
zu ihm unbegrenzt Material sammeln, ohne je fertig zu werden, weil nichts
festlegt, wann man genug weiß. Erst die Forschungsfrage schneidet aus dem
Gebiet ein Stück heraus, das sich bearbeiten lässt. \textcite{booth2016craft}
schlagen dafür eine einfache Probe vor: Man solle den Satz \enquote{Ich
untersuche X, weil ich herausfinden will, Y, damit andere Z verstehen}
vollständig ausfüllen können. X ist das Thema, Y die Frage, Z die Antwort
auf das \enquote{Na und?}. Wer bei Y oder Z ins Stocken kommt, hat noch
keine Forschungsfrage.}
\begin{description}
\item[{\rot\bf Vier Begriffe, die oft verwechselt werden}]~\\[-\lblskip]
\begin{center}
\small
\begin{tabular}{p{0.22\textwidth}p{0.66\textwidth}}
\toprule
\textbf{Thema} & das Gebiet: \enquote{Zauberkunst im Physikunterricht} \\
\textbf{Erkenntnis\-interesse} & warum mich das angeht: \enquote{Ich möchte
  wissen, ob Staunen beim Lernen hilft.} \\
\textbf{Forschungsfrage} & was genau geklärt werden soll: \enquote{Behalten
  Schüler:innen der 8.~Klasse ein physikalisches Prinzip besser, wenn es mit
  einem Zaubereffekt eingeführt wird?} \\
\textbf{Hypothese} & die vermutete, prüfbare Antwort: \enquote{Sie behalten
  es besser.} ($\to$~LE~5) \\
\bottomrule
\end{tabular}
\end{center}
\os{\newpage}

\item[{\rot\bf Die Probe nach Booth u.\,a.}]~\\[-\lblskip]
\bi
  \ite Ich untersuche \textbf{X} \dots
  \ite weil ich herausfinden will, \textbf{Y} \dots
  \ite damit andere \textbf{Z} verstehen.
  \ite wer Y nicht ausfüllen kann, hat erst ein Thema
  \ite wer Z nicht ausfüllen kann, hat noch keine Relevanz
    \parencite{booth2016craft}
\ei
\os{\newpage}

\item[{\rot\bf Arbeitsdefinition}]~\\
\defn{forschungsfrage}{Eine \textbf{Forschungsfrage} legt fest, welcher
Sachverhalt bei welchem Gegenstand unter welchen Bedingungen geklärt werden
soll, und zwar so genau, dass man angeben kann, welche Daten sie
beantworten würden.}
\bi
  \ite sie bestimmt, wonach man sucht und was man weglassen darf
  \ite sie ist kleiner als das Thema und größer als eine Ja-Nein-Auskunft
\ei
\os{\newpage}

\item[{\rot\bf Arten von Forschungsfragen}]~\\[-\lblskip]
\bi
  \ite \textbf{beschreibend:} Wie verbreitet ist X? Wie sieht X aus?
  \ite \textbf{erklärend:} Warum tritt X auf? Was bewirkt X?
  \ite \textbf{vergleichend:} Unterscheiden sich A und B in X?
  \ite \textbf{verstehend:} Wie erleben Betroffene X?
  \ite \textbf{gestaltend:} Wie muss Y aussehen, damit X erreicht wird?
\ei
\nsl{Die Fragetypen legen bereits eine Methodenrichtung nahe
\parencite{doering2016forschungsmethoden}. Beschreibende und vergleichende
Fragen führen häufig zu Befragungen oder Beobachtungen mit Zählungen,
erklärende Fragen zu Experimenten oder Längsschnittstudien, verstehende zu
qualitativen Interviews ($\to$~LE~9). Wer eine Frage formuliert, trifft
also schon eine Vorentscheidung über das Design ($\to$~LE~6).}
\os{\newpage}
\end{description}

%==============================================================================
\section{Neugier als Informationslücke}
%==============================================================================
\nsl{Warum lohnt sich der Aufwand, aus einer vagen eine scharfe Frage zu
machen? Das Buch gibt in Kapitel~4 eine kognitionspsychologische Antwort.
George Loewenstein beschreibt Neugier als einen Zustand der Entbehrung,
vergleichbar mit Hunger: Sie entsteht, wenn wir eine Lücke zwischen dem
bemerken, was wir wissen, und dem, was wir wissen wollen
\parencite{loewenstein1994}. Man muss die Lücke allerdings \emph{bemerken}.
Wer gar nichts über ein Gebiet weiß, sieht keine Lücke; wer schon alles zu
wissen glaubt, auch nicht. Neugier ist dort am stärksten, wo man schon
etwas weiß und genau sieht, was fehlt.

Bereits \textcite{berlyne1954} hatte zwischen \emph{spezifischer} und
\emph{diversiver} Neugier unterschieden: dem gerichteten Wunsch nach einer
bestimmten fehlenden Information und dem diffusen Bedürfnis nach
Abwechslung. Ein Zaubertrick erzeugt nach der Überlegung im Buch
(Kap.~7d) die erste Form, denn er stellt genau eine Frage: \enquote{Wie
kann das sein?} Eine gute Forschungsfrage leistet dasselbe auf
kontrollierte Weise. Sie macht aus diffusem Interesse eine scharf umrissene
Lücke.}
\begin{description}
\item[{\rot\bf Neugier entsteht an der Lücke}]~\\[-\lblskip]
\begin{center}
\resizebox{0.85\textwidth}{!}{%
\begin{tikzpicture}[font=\sffamily,>=Stealth]
  \fill[rwuviolet!15] (0,0) rectangle (5,1.2);
  \node at (2.5,0.6) {was ich weiß};
  \fill[rwucyan40!50] (5,0) rectangle (7.5,1.2);
  \draw[rwucyan,very thick] (5,0) rectangle (7.5,1.2);
  \node at (6.25,0.6) {\textbf{Lücke}};
  \draw[rwugray] (0,0) rectangle (10,1.2);
  \node[rwugray,font=\sffamily\small] at (8.75,0.6) {unbemerkt};
  \draw[<->,thick] (0,1.5) -- node[above]{was ich wissen will} (7.5,1.5);
  \node[align=center,text width=9.5cm] at (5,-0.8) {Neugier ist am größten,
    wenn die Lücke \emph{sichtbar} und \emph{begrenzt} ist};
\end{tikzpicture}}
\end{center}
\bi
  \ite Neugier als Entbehrungszustand, wie Hunger \parencite{loewenstein1994}
\ei
\os{\newpage}

\item[{\rot\bf Spezifisch statt diffus}]~\\[-\lblskip]
\bi
  \ite \textbf{diversive Neugier:} \enquote{Irgendwas über Social Media
    wäre interessant.}
  \ite \textbf{spezifische Neugier:} \enquote{Warum scrolle ich abends länger,
    als ich vorhatte?} \parencite{berlyne1954}
  \ite eine Forschungsfrage verwandelt die erste in die zweite
  \ite Faustregel: Wer nicht sagen kann, wann er die Antwort hätte, hat
    die Lücke noch nicht umrissen
\ei
\os{\newpage}

\item[{\rot\bf Das wohldefinierte Problem}]~\\[-\lblskip]
\bi
  \ite Einstein fragte nicht \enquote{Wie funktioniert das Universum?}
  \ite sondern: Warum beschreibt die Elektrodynamik bewegten Magneten und
    bewegten Leiter verschieden? \parencite{einstein1905}
  \ite ein scharfes Problem verkleinert den Suchraum
    \parencite{epstein2026insidebox}
  \ite \enquote{Lass eine Münze verschwinden} ist unlösbar; \enquote{unter
    diesen Bedingungen, für dieses Publikum} wird lösbar
\ei
\nsl{Das Buch greift in Kapitel~7d eine These von
\textcite{epstein2026insidebox} auf: Wissenschaftliche Durchbrüche entstehen
oft nicht aus freier Spekulation, sondern aus der hartnäckigen Arbeit an
einem eng gestellten Problem. Das Lehrstück ist der Anfang der Arbeit von
1905, in der Einstein von einer Asymmetrie ausgeht, die \enquote{den
Phänomenen nicht anzuhaften scheint} \parencite{einstein1905}: Für die
Theorie macht es einen Unterschied, ob sich der Magnet oder der Leiter
bewegt, für den gemessenen Strom nicht. Ein solches Problem wirkt wie eine
Begrenzung, die Energie bündelt. Für Studierende ist das entlastend. Eine
kleine, genaue Frage zeugt nicht von fehlendem Ehrgeiz; sie ist die
Voraussetzung dafür, überhaupt eine Antwort zu finden.}
\os{\newpage}
\end{description}

%==============================================================================
\section{Kriterien guter Forschungsfragen}
%==============================================================================
\nsl{Für die Prüfung einer Forschungsfrage gibt es mehrere Kriterienlisten,
die sich stark überschneiden. Das Curriculum nennt vier Grundkriterien:
Relevanz, Präzision, Beantwortbarkeit und Originalität. Zwei Merkhilfen
aus anderen Fachkulturen sind verbreitet. Das SMART-Schema stammt aus dem
Management und beschreibt dort gut formulierte Ziele
\parencite{doran1981smart}; in der Hochschullehre wird es gern auf
Forschungsvorhaben übertragen. Die FINER-Kriterien kommen aus der
klinischen Forschung \parencite{hulley2013designing} und betonen
zusätzlich Durchführbarkeit und Ethik. Keine der Listen ist ein Algorithmus.
Es sind Prüffragen, die man an einen Entwurf richtet.}
\begin{description}
\item[{\rot\bf Vier Grundkriterien}]~\\[-\lblskip]
\bi
  \ite \textbf{Relevanz:} Wen interessiert die Antwort, und warum?
  \ite \textbf{Präzision:} Sind Gegenstand, Personen, Ort und Zeit
    festgelegt?
  \ite \textbf{Beantwortbarkeit:} Gibt es Daten, die die Frage entscheiden
    könnten?
  \ite \textbf{Originalität:} Ist die Antwort nicht schon bekannt?
\ei
\os{\newpage}

\item[{\rot\bf SMART auf Forschungsfragen übertragen}]~\\[-\lblskip]
\begin{center}
\small
\begin{tabular}{lp{0.62\textwidth}}
\toprule
\textbf{S} spezifisch & ein Gegenstand, eine Population, ein Kontext \\
\textbf{M} messbar & es ist klar, welche Größe erhoben wird \\
\textbf{A} erreichbar & Zugang zu Daten, Personen, Material \\
\textbf{R} relevant & die Antwort verändert etwas (Wissen, Praxis) \\
\textbf{T} terminiert & im Semester (oder in der Arbeit) zu schaffen \\
\bottomrule
\end{tabular}
\end{center}
\bi
  \ite ursprünglich ein Schema für Managementziele \parencite{doran1981smart}
\ei
\os{\newpage}

\item[{\rot\bf FINER: die Variante aus der Medizin}]~\\[-\lblskip]
\bi
  \ite \textbf{F}easible: durchführbar (Zeit, Geld, Fallzahl)
  \ite \textbf{I}nteresting: für Sie und für andere interessant
  \ite \textbf{N}ovel: bestätigt, widerlegt oder erweitert Bekanntes
  \ite \textbf{E}thical: vertretbar für alle Beteiligten ($\to$~LE~13)
  \ite \textbf{R}elevant: folgenreich für Wissenschaft oder Praxis
    \parencite{hulley2013designing}
\ei
\os{\newpage}

\item[{\rot\bf Typische Schwächen}]~\\[-\lblskip]
\bi
  \ite \textbf{zu weit:} \enquote{Wie wirkt Digitalisierung auf die
    Bildung?}
  \ite \textbf{zu eng, trivial:} \enquote{Wie viele Hörsäle hat der
    Campus?}
  \ite \textbf{normativ:} \enquote{Sollte man Handys an Schulen verbieten?}
  \ite \textbf{nicht entscheidbar:} \enquote{Hat ein Zaubertrick eine
    Seele?}
  \ite \textbf{Antwort eingebaut:} \enquote{Warum schadet TikTok der
    Konzentration?}
\ei
\nsl{Normative Fragen sind nicht unwissenschaftlich, aber sie lassen sich
nicht durch Daten allein entscheiden. Man kann sie in empirische Teilfragen
zerlegen (Was ändert sich an Schulen nach einem Handyverbot?) und die
Bewertung offen ausweisen. Die letzte Schwäche führt zum Staunmoment
zurück: Die Frage setzt voraus, was sie eigentlich prüfen müsste.}
\os{\newpage}
\end{description}

%==============================================================================
\section{Implizite Prämissen}
%==============================================================================
\nsl{Der Satz über den Becher war wahr, aber er ließ die Zuhörenden eine
Voraussetzung ergänzen, die nicht dastand. Dasselbe geschieht in Fragen.
Die Sprachphilosophie spricht von \emph{Präsuppositionen}: Annahmen, die
eine Äußerung voraussetzt, ohne sie zu behaupten. Die Frage \enquote{Warum
schadet TikTok der Konzentration?} setzt voraus, dass es einen Schaden
gibt; wer sie beantwortet, hat die Voraussetzung schon akzeptiert. In der
Rhetorik heißt das die \emph{suggestive} oder \emph{komplexe Frage}.

Wie wenig harmlos solche Voraussetzungen sind, zeigt ein klassisches
Experiment. \textcite{loftus1974reconstruction} ließen Versuchspersonen
Filme von Autounfällen sehen und fragten danach, wie schnell die Wagen
gefahren seien, als sie \enquote{zusammenstießen} (\emph{hit}) oder
\enquote{ineinanderkrachten} (\emph{smashed}). Das stärkere Verb führte zu
höheren Schätzungen, und eine Woche später berichteten mehr Personen von
zerbrochenem Glas, das im Film gar nicht vorkam. Die Frage hat die Antwort
mitgeformt. Für das wissenschaftliche Arbeiten folgt daraus: Eine
Forschungsfrage muss so gebaut sein, dass auch die unerwartete Antwort
möglich bleibt.}
\begin{description}
\item[{\rot\bf Was eine Frage voraussetzt}]~\\[-\lblskip]
\bi
  \ite \textbf{Existenz:} \enquote{Wie stark wirkt der Mondzyklus auf den
    Schlaf?} setzt voraus, dass er wirkt
  \ite \textbf{Richtung:} \enquote{Warum macht Homeoffice produktiver?}
  \ite \textbf{Messbarkeit:} \enquote{Wie kreativ sind Erstsemester?}, als
    sei Kreativität eine Zahl
  \ite \textbf{Gruppe:} \enquote{die Studierenden}: welche, wo, wann?
\ei
\os{\newpage}

\item[{\rot\bf Fragen formen Antworten}]~\\[-\lblskip]
\bi
  \ite Unfallfilm, dann die Frage: Wie schnell waren die Autos, als sie
    \enquote{zusammenstießen} oder \enquote{ineinanderkrachten}?
  \ite stärkeres Verb $\to$ höhere Schätzungen
  \ite eine Woche später: mehr Erinnerungen an Glassplitter, die es nicht
    gab \parencite{loftus1974reconstruction}
  \ite die Formulierung ist Teil der Messung
\ei
\os{\newpage}

\item[{\rot\bf Prämissen freilegen: drei Prüffragen}]~\\[-\lblskip]
\bi
  \ite Was muss wahr sein, damit die Frage überhaupt Sinn ergibt?
  \ite Welche Antwort ist \emph{ausgeschlossen}, so wie die Frage gestellt
    ist?
  \ite Welches Wort würde jemand anderes anders verstehen?
\ei
\nsl{Die dritte Prüffrage ist die Brücke zur Operationalisierung. Im
Staunmoment war es das Wort \enquote{etwas}: Die Gruppe hat es als
\enquote{eine Art von Ding} gelesen, gemeint war \enquote{ein einzelnes
Exemplar}. In Forschungsfragen sind es Wörter wie \enquote{Erfolg},
\enquote{Stress}, \enquote{Lernen} oder \enquote{besser}, die jeder
anders füllt.}
\os{\newpage}
\end{description}

%==============================================================================
\section{Operationalisierung: vom Begriff zur Messung}
%==============================================================================
\nsl{Viele Forschungsfragen handeln von Dingen, die man nicht direkt sehen
kann: Neugier, Lernerfolg, Vertrauen, Staunen. Solche Begriffe heißen in
der Methodenlehre \emph{Konstrukte}. Um sie zu untersuchen, braucht man
beobachtbare \emph{Indikatoren} und eine genaue \emph{Messvorschrift}. Den
Weg vom Konstrukt zur Messung nennt man Operationalisierung
\parencite{doering2016forschungsmethoden}. Der Begriff geht auf den
Physiker Percy Bridgman zurück, für den ein Begriff durch die Operationen
bestimmt war, mit denen man ihn misst \parencite{bridgman1927logic}. In
dieser strengen Form hat sich die Idee nicht gehalten, denn dann wäre jede
neue Messmethode ein neuer Begriff. Die gemäßigte Fassung gehört aber zum
Grundbestand jeder empirischen Arbeit: Wer nicht angibt, wie er misst,
kann weder zeigen, dass er recht hat, noch dass er unrecht hat.

Operationalisierung ist Verlust und Gewinn zugleich. Ein Test über
Fachbegriffe erfasst nicht \enquote{das Lernen}, nur einen Ausschnitt
davon. Dafür kann man ihn wiederholen, vergleichen und anfechten. Die Kunst
liegt darin, den Kern der staunenswerten Frage zu bewahren und ihn
trotzdem messbar zu machen.}
\begin{description}
\item[{\rot\bf Vom Konstrukt zur Messvorschrift}]~\\[-\lblskip]
\begin{center}
\resizebox{0.95\textwidth}{!}{%
\begin{tikzpicture}[font=\sffamily,>=Stealth,
  st/.style={draw=rwuviolet,thick,fill=rwuviolet!8,rounded corners=2pt,
    minimum height=20mm,minimum width=37mm,align=center,text width=33mm}]
  \node[st] (k) at (0,0) {\textbf{Konstrukt}\\\enquote{Lernerfolg}};
  \node[st] (d) at (4.2,0) {\textbf{Definition}\\Behalten eines\\Prinzips};
  \node[st] (i) at (8.4,0) {\textbf{Indikator}\\Testaufgaben\\zum Prinzip};
  \node[st,draw=rwucyan,fill=rwucyan40!40] (m) at (12.6,0) {\textbf{Messung}\\Punkte im Test,\\eine Woche später};
  \draw[->,very thick] (k) -- (d);
  \draw[->,very thick] (d) -- (i);
  \draw[->,very thick] (i) -- (m);
  \draw[rwugray,thick,decorate,decoration={brace,mirror,amplitude=6pt}]
    (-1.8,-2.0) -- node[below=7pt]{abstrakt, umstritten}(6.0,-2.0);
  \draw[rwugray,thick,decorate,decoration={brace,mirror,amplitude=6pt}]
    (6.6,-2.0) -- node[below=7pt]{konkret, wiederholbar, angreifbar}(14.4,-2.0);
\end{tikzpicture}}
\end{center}
\os{\newpage}

\item[{\rot\bf Was eine Messvorschrift festlegt}]~\\
\defn{operationalisierung}{\textbf{Operationalisierung} heißt, für einen
Begriff anzugeben, durch welche beobachtbaren Größen er erfasst wird und
nach welcher Regel aus Beobachtungen ein Messwert entsteht.}
\bi
  \ite \emph{was} beobachtet wird (Indikator)
  \ite \emph{wie} und \emph{wann} (Instrument, Zeitpunkt)
  \ite \emph{wie} daraus ein Wert wird (Auswertungsregel)
\ei
\os{\newpage}

\item[{\rot\bf Gewinn und Verlust}]~\\[-\lblskip]
\bi
  \ite gewonnen: Wiederholbarkeit, Vergleichbarkeit, Angreifbarkeit
  \ite verloren: alles am Begriff, was der Indikator nicht erfasst
  \ite darum wo möglich mehrere Indikatoren ($\to$~Triangulation, LE~6)
  \ite Prüffrage: Würde eine Kollegin bei dieser Messung dasselbe unter dem
    Begriff verstehen?
\ei
\nsl{Die Prüffrage führt auf das Gütekriterium der Validität, das in LE~6
ausführlich behandelt wird. Für heute genügt die Einsicht, dass jede
Operationalisierung eine Entscheidung ist, die man begründen und offenlegen
muss.}
\os{\newpage}

\item[{\rot\bf Das geheime Objekt, operationalisiert}]~\\[-\lblskip]
\bi
  \ite \enquote{noch kein Mensch gesehen}: Was zählt als \enquote{sehen}?
    Die Schale? Ein Röntgenbild?
  \ite \enquote{etwas}: Exemplar oder Art?
  \ite \enquote{nach dieser Stunde}: Wie würde man prüfen, dass niemand es
    mehr sieht?
  \ite erst mit diesen Festlegungen wird der Satz prüfbar
\ei
\os{\newpage}
\end{description}

%==============================================================================
\section{Eingrenzen und die Forschungslücke finden}
%==============================================================================
\nsl{Eine Frage ist erst dann eine Forschungsfrage, wenn ihre Antwort
nicht schon vorliegt. Die \emph{Forschungslücke} ist der Teil des
Gegenstands, zu dem es noch keine belastbare Antwort gibt: weil eine
Gruppe noch nicht untersucht wurde, weil Befunde einander widersprechen,
weil eine Methode noch nicht angewandt wurde oder weil ein bekannter Befund
noch nie wiederholt wurde. Ob eine Lücke besteht, zeigt erst die
Literaturrecherche ($\to$~LE~3). Im Bachelorstudium muss die Lücke nicht
groß sein. Eine Replikation unter anderen Bedingungen oder die Übertragung
eines bekannten Befunds auf einen neuen Kontext ist eine legitime Lücke.}
\begin{description}
\item[{\rot\bf Der Trichter}]~\\[-\lblskip]
\begin{center}
\resizebox{0.8\textwidth}{!}{%
\begin{tikzpicture}[font=\sffamily]
  \fill[rwuviolet!12] (-2,4) -- (12,4) -- (8,0) -- (2,0) -- cycle;
  \draw[rwuviolet,thick] (-2,4) -- (2,0) -- (8,0) -- (12,4);
  \node at (5,3.5) {\textbf{Thema:} Zauberkunst und Lernen};
  \node at (5,2.6) {Fach: Physikunterricht};
  \node at (5,1.8) {Population: 8. Klasse};
  \node at (5,0.9) {Größe: Behalten nach 1 Woche};
  \node[fill=rwucyan40!60,draw=rwucyan,rounded corners=2pt,align=center] at (5,-0.8)
    {\textbf{Forschungsfrage:} Behalten \dots\ besser, wenn \dots ?};
\end{tikzpicture}}
\end{center}
\os{\newpage}

\item[{\rot\bf Woran man eine Forschungslücke erkennt}]~\\[-\lblskip]
\bi
  \ite Gruppe, Ort oder Zeitraum wurden noch nicht untersucht
  \ite Befunde widersprechen einander
  \ite eine Methode wurde noch nicht angewandt
  \ite ein Befund wurde noch nie wiederholt (Replikation)
  \ite die Belege kommen aus der Recherche, nicht aus dem Bauchgefühl
\ei
\os{\newpage}

\item[{\rot\bf Staunenswerte und beantwortbare Fragen}]~\\[-\lblskip]
\begin{center}
\small
\begin{tabular}{p{0.3\textwidth}p{0.58\textwidth}}
\toprule
\textbf{staunenswert} & \textbf{beantwortbar} \\
\midrule
Warum staunen wir? & Stellen Studierende nach einem Zaubereffekt mehr
  Nachfragen als nach einem gleich langen Experiment ohne Überraschung? \\
Macht Musik klüger? & Lösen Studierende nach 10 Min.\ Musik mehr
  Aufgaben eines Konzentrationstests als nach 10 Min.\ Stille? \\
Was ist ein gutes Team? & Welche Merkmale nennen Mitglieder von
  Projektteams im 4.~Semester als förderlich? \\
\bottomrule
\end{tabular}
\end{center}
\bi
  \ite die linke Spalte ist der Motor, die rechte der Plan
\ei
\os{\newpage}
\end{description}

\begin{merkbox}
Eine Forschungsfrage ist eine Informationslücke mit Rand: Sie sagt, was
man wissen will, welche Voraussetzungen man dabei macht und welche
Beobachtung die Frage beantworten würde. Wer die Prämissen seiner Frage
nicht kennt, findet am Ende die Antwort, die er hineingelegt hat.
\end{merkbox}
\os{\newpage}

%==============================================================================
\section{Übung in der Stunde}
%==============================================================================
\begin{description}
\item[{\rot\bf Werkstatt: Fragen sezieren (Zweiergruppen, 30 Min.)}]~\\[-\lblskip]
\bi
  \ite Aufwärmen: zwei Fragen aus der Liste auf dem Übungsblatt. Welche
    Prämisse steckt darin?
  \ite Dann Ihre eigene Frage: Was wird als selbstverständlich
    vorausgesetzt?
  \ite Welche Wörter müssten Sie operationalisieren? Was würden Sie
    messen?
  \ite Schärfen Sie die Frage nach SMART; Ihr Gegenüber prüft
\ei
\os{\newpage}

\item[{\rot\bf Sicherung (15 Min.)}]~\\[-\lblskip]
\bi
  \ite Schreiben Sie Ihre Forschungsfrage auf eine Karteikarte
  \ite Plenum: Ist sie beantwortbar? Was wäre die Datengrundlage?
  \ite Welche Antwort würde Sie überraschen?
\ei
\os{\newpage}
\end{description}

\begin{lehrhinweis}
In der Werkstatt geht die Lehrperson von Paar zu Paar und stellt immer
dieselben drei Fragen: \enquote{Was muss wahr sein, damit Ihre Frage Sinn
ergibt?}, \enquote{Welche Antwort würde Sie überraschen?} und
\enquote{Woran würden Sie die Antwort erkennen?} Häufige Muster bei
mitgebrachten Fragen sind normative Fragen (\enquote{Sollte man \dots?}),
Warum-Fragen mit eingebauter Richtung und Fragen ohne Population. Die
Studierenden sollen ihre Frage nicht aufgeben, sondern umbauen. Aus
\enquote{Warum macht Sport glücklicher?} wird zum Beispiel \enquote{Hängt
bei Studierenden der RWU die Zahl der Sporteinheiten pro Woche mit der
selbstberichteten Lebenszufriedenheit zusammen?} Am Ende der Sicherung
darauf hinweisen, dass eine Zusammenhangsfrage noch keine Ursache klärt
(Vorgriff auf LE~6). Für die Plenumsrunde drei bis vier Karten auswählen,
die typische Schwächen zeigen, und die Verfasser:innen vorher fragen, ob
sie einverstanden sind.
\end{lehrhinweis}

%==============================================================================
\section{Lernziele und Aufgaben}
%==============================================================================
\begin{description}
\item[{\rot\bf Nach dieser Einheit können Sie}]~\\[-\lblskip]
\bi
  \ite Thema, Erkenntnisinteresse, Forschungsfrage und Hypothese
    unterscheiden
  \ite implizite Prämissen in Fragestellungen erkennen und benennen
  \ite einen Begriff operationalisieren
  \ite eine eigene Forschungsfrage entwickeln und nach SMART prüfen
\ei
\os{\newpage}

\item[{\rot\bf Bis zur nächsten Woche}]~\\[-\lblskip]
\bi
  \ite \textbf{Forschungsfrage schriftlich:} Formulieren Sie Ihre Frage
    und begründen Sie in drei Sätzen, warum sie beantwortbar und nicht
    trivial ist
  \ite \textbf{Probe nach Booth:} \enquote{Ich untersuche X, weil \dots,
    damit \dots} vollständig ausfüllen
  \ite \textbf{Tagebuch des Staunens:} Formulieren Sie Ihr Staunerlebnis
    dieser Woche als Frage. Welche Prämisse steckt darin?
  \ite Bringen Sie die Frage zu LE~3 mit: Mit ihr starten Sie die
    Recherche
\ei
\os{\newpage}
\end{description}

\handoutinput{handouts/LE02}
